\chapter{Discussion, Limitations, and Future Work}
\label{ch:discussion}

The preceding chapters have made a sustained argument that standing must
be earned through admission rather than asserted by prose, and
Chapter~\ref{ch:evaluation} closed by enumerating the eight closed
failure classes that would falsify the release if any were inhabited.
This chapter turns that same scrutiny on the release's own boundaries:
what the proven/stated discipline does and does not claim, what kernel
admission does and does not check, what this thesis explicitly excludes,
and --- in the Future Work section --- what remains open, including a
frank account of a toolchain limitation specific to the sandbox this
thesis was written in, kept carefully separate from any claim about the
certified release itself.

\section{Proven versus Stated}
\label{sec:disc-proven-stated}

\texttt{paper/main.tex}'s Limitations and Standing section states the
library's discipline plainly: the library's standing is exactly the
manifest's proven/stated split, and prose adds nothing to it. A
\emph{stated} declaration, in this vocabulary, is a formal statement
whose proof obligation is open: it compiles under the pinned toolchain,
it is citable as a definition of the property it names, and it carries no
more authority than that. Chapter~\ref{ch:procint} exercised this
distinction on the \TotalDecls-declaration catalog, and it is worth
restating the current split precisely, without smoothing the two
categories together as this thesis's discipline forbids: \KernelTheorems\
declarations are kernel-admitted, axiom-audited proven theorems, and
\StatedCount\ remain formally stated with an open obligation.

Among those \StatedCount\ stated declarations, two categories are worth
distinguishing again here, because collapsing them is exactly the kind of
mistake the proven/stated discipline exists to prevent. The first is a
single declaration, \texttt{ProcInt.BranchingProcess.isUnfoldingOf\_statement}
(\texttt{packs/lean-math-pack/fragments/petriext.ttl}), the unfolding-correctness
statement tracked in \texttt{research/wfnet/obligations.toml} as
obligation order 3, \texttt{name = "unfolding\_correctness"},
\texttt{status = "STATED"} --- this is the one declaration in the current
catalog that represents genuine open proof work, in the sense that
closing it would add a new kernel-admitted theorem rather than merely
documenting an intentional boundary. The second category is the five
WF-net countermodel declarations in
\texttt{packs/lean-math-pack/fragments/workflow\_countermodel.ttl} ---
\texttt{crownCounter\_reaches\_mid}, \texttt{crownCounter\_reaches\_final},
\texttt{crownCounter\_sound}, \texttt{crownCounter\_not\_bounded}, and
\texttt{WfNet.infinite\_transition\_countermodel\_sound\_not\_bounded} ---
which are \emph{deliberately and permanently} stated. These five
declarations do not represent unfinished work; they witness why the
crown-jewel equivalence requires its \texttt{[Finite~T]} hypothesis in
the first place, by exhibiting a net that is sound and not bounded once
that hypothesis is dropped. A guard, \texttt{countermodel\_not\_promoted},
is engineered specifically to refuse (\texttt{COUNTERMODEL\_PROMOTION\_REFUSED})
any attempt to promote these five to proven standing, precisely because
promoting a deliberate negative witness to proven status would corrupt
the very fact it is meant to record. \texttt{just status} in this
sandbox's environment reports this family under its own status key,
\texttt{WFNET\_INFINITE\_TRANSITION\_COUNTERMODEL=STATED}, distinct from
the ordinary stated/proven split, for exactly this reason: it is a
status that is correct and intended to remain STATED, not a gap awaiting
closure alongside the unfolding-correctness obligation.

\section{Crown-Jewel Scope}
\label{sec:disc-crown-jewel}

The crown-jewel theorem --- the workflow-net soundness equivalence
(van der Aalst 1997, Lemma 8)~\cite{vanderaalst1997}, characterizing
soundness as equivalent to short-circuit liveness and boundedness under
\texttt{[Finite~T]} --- carries \CrownJewelStatus\ status at the current
release. \texttt{paper/main.tex} is careful to scope this claim per
direction: if only one direction of the equivalence were admitted at a
given release, the claim would be that direction and nothing more, with
the converse remaining a stated obligation in the manifest rather than an
implicit part of the headline result. This chapter reproduces that
caution rather than restating a stronger claim than the manifest
supports; the crown-jewel status fragment rendered in
Chapter~\ref{ch:evaluation} and Chapter~\ref{ch:procint} is the authority
on which direction, or directions, the current release actually admits.

\section{What Admission Does Not Check}
\label{sec:disc-admission-limits}

Kernel admission together with the axiom audit establishes a specific
and narrow thing: that an admitted term proves its stated theorem under
the pinned, authorized axiom set. It does not establish that the
catalog-curated statement is a faithful rendering of the informal
mathematical canon it claims to formalize. That correspondence --- does
this Lean proposition actually say what van der Aalst's Lemma~8 says, or
what the process-mining literature means by ``sound'' --- rests on
citation-annotated statement curation reviewed against primary sources,
and is exactly the kind of residual risk the semantic audits of Meek et
al.~\cite{meek2026formalizing} target: compilation success can coexist
with a silently weakened hypothesis or an incompletely stated multi-part
theorem, and no amount of kernel-side checking distinguishes a faithful
formalization from a subtly narrower one that still type-checks.
Chapter~\ref{ch:background} introduced this diagnosis; this chapter
returns to it as a standing limitation rather than a risk this pipeline
claims to have retired. Fidelity of statements is, in this repository's
own words, a review obligation, not a theorem.

A second limit concerns novelty. Any ``first mechanization'' claim made
elsewhere in this thesis or in \texttt{paper/main.tex} is scoped to a
search of the Lean (Mathlib/CSLib), Coq, and Isabelle (AFP) ecosystems
conducted at run time; it is a negative search result at a point in time,
not a proof of absence, and later work in any of those ecosystems could
in principle supersede it without this thesis's claim having been false
when made.

A third limit is that standing is not to be inferred from this document's
prose at all. At release time, the crown-jewel rail, the countermodel
rail, the correspondence rail, the product/DX rail, and the paper rail
each carry independent manifest status; the paper, and by extension this
thesis, is a publication surface reporting that status, not itself a
source of truth. Where paper prose and manifest status diverge, the
manifest and its receipts govern, and this chapter treats that as a
standing methodological commitment rather than a hypothetical.

\section{Exclusions}
\label{sec:disc-exclusions}

\texttt{paper/main.tex} enumerates, in tabular form, a set of claims this
work explicitly does not make, and this thesis inherits that enumeration
rather than relaxing it: that LLMs prove mathematics autonomously (they
are treated throughout as untrusted candidate producers, per
Chapter~\ref{ch:mfact} and Chapter~\ref{ch:ai-development}); that
\textsf{procint} formalizes all of process intelligence (the v1 boundary
is manifest-scoped, per Chapter~\ref{ch:procint}); that kernel admission
proves informal intent (statement fidelity is a review obligation, per
Section~\ref{sec:disc-admission-limits} above); that Rice's theorem is
violated (arbitrary semantic correctness is not decided by anything in
this pipeline); that external opinion enters the standing calculus (only
receipted artifacts are parsed by the certification gate); that criticism
outside the declared boundary is refuted by the empty-objection-surface
claim of Chapter~\ref{ch:evaluation} (that claim is scoped strictly to
the declared boundary and says nothing about objections outside it); that
\textsf{ggen} output has standing merely by having been generated (only
admitted output has standing, per the projection-versus-admission
distinction of Chapter~\ref{ch:mfact}); and that a stated theorem is
proven (stated means, without qualification, that the obligation remains
open).

\section{Future Work}
\label{sec:disc-future-work}

Four items make up this thesis's honest account of what remains open,
ranging from a specific closable proof obligation to a limitation of this
particular writing environment that this chapter is careful not to
conflate with the certified release's own standing.

\paragraph{(a) Closing the unfolding-correctness obligation.}
The single genuinely open item in the current \TotalDecls-declaration
catalog, \texttt{ProcInt.BranchingProcess.isUnfoldingOf\_statement}, is a
concrete, scoped target for future proof work: its statement, as recorded
in \texttt{packs/lean-math-pack/fragments/petriext.ttl}, characterizes
when a branching process is a legal unfolding of a net~$N$ by conjoining
acyclicity of the occurrence net, absence of backward conflict, and two
conditions relating the occurrence net's pre/post-condition structure to
$N$'s transitions on singleton arcs. The declaration's own doc-comment
attributes the underlying homomorphism condition it specializes to
Engelfriet 1991, Definition 3.1; this thesis reports that attribution
exactly as the source states it, in prose rather than as a citation
key, because no corresponding entry exists in \texttt{paper/refs.bib} and
this thesis's discipline forbids inventing one. Closing this obligation
means discharging its stated proposition under the same kernel-admission
and axiom-audit gate every other \KernelTheorems-count theorem in the
catalog passed through --- not weakening the statement to make it easier,
and not asserting proven status without the mechanical check.

\paragraph{(b) Un-run docs and post-release lanes.}
\texttt{just next}, run read-only in this sandbox, lists two lanes that
have not yet been attempted in this environment: the publication packet
has not been manufactured (actionable via \texttt{just manufacture-post-release}),
and the docs lane has not been attempted (actionable via \texttt{just docs}).
Neither omission changes the certified release's proven/stated split or
gate status; they are lanes downstream of certification, not inputs to
it, and their non-execution here is reported as an open item rather than
elided.

\paragraph{(c) External-actuation packets.}
\texttt{just status} in this sandbox reports \texttt{PUBLICATION\_ACTUATION=PENDING\_EXTERNAL\_ACTUATION},
with the arXiv submission packet (\texttt{ARXIV\_PACKET}) and the
GitHub-push packet (\texttt{GITHUB\_PUSH\_PACKET}) both currently
\texttt{BLOCKED}, while the GitHub-release packet
(\texttt{GITHUB\_RELEASE\_PACKET}) reports \texttt{ALIVE}. In this
repository's status taxonomy, \texttt{PENDING\_EXTERNAL\_ACTUATION}
denotes a packet that is itself complete but whose remaining step ---
submitting to arXiv, pushing to a remote the agent does not have
standing to push to --- is the user's action, not a further manufacturing
step this pipeline can discharge on its own. Reporting these packets as
blocked or pending is, consistent with this thesis's discipline
throughout, preferable to describing the release as ``published'' before
that external step has actually occurred.

\paragraph{(d) A sandbox toolchain limitation, scoped to this environment.}
This chapter closes with an honest report of a limitation particular to
the sandbox this thesis was written in, stated carefully so that it is
not read as a claim about the certified \ReleaseTag\ release itself.
Checked in this sandbox on 2026-07-21: \texttt{lake} and \texttt{elan}
are not on \texttt{PATH}, \texttt{ggen} is not on \texttt{PATH}, and
\texttt{procint}'s and \textsf{mfact}'s pinned toolchain
(\texttt{leanprover/lean4:v4.31.0}, per \texttt{lean-toolchain}) is not
installed in this sandbox's \texttt{elan} toolchain list, which is empty.
\texttt{just doctor}, itself a read-only diagnostic that does not dirty
the tree, confirms each of these and additionally reports that the
\texttt{lean-math-pack}, \texttt{quadrature-pack}, and
\texttt{post-release-pack} sources are absent, because they live at a
\texttt{/Users/sac/mfact} path this sandbox --- rooted at
\texttt{/home/user/mfact} --- does not have. A single bounded
installation attempt, \texttt{just install}, was run per this
repository's actuation discipline (Chapter~\ref{ch:ai-development}); it
proved to be a checker rather than an installer, printing guidance
(\texttt{elan: not found --- install with: curl https://elan.dev/install.sh -sSf | sh};
an analogous \texttt{ggen} instruction pointing at \texttt{just install-ggen}
from a praxis checkout) without itself performing a network fetch, and no
further installation was attempted, consistent with this repository's
instruction not to bypass its own guarded actuation paths. Toolchain
availability in this sandbox was unchanged, unavailable both before and
after this bounded attempt.

None of this bears on the certified release's own standing. \texttt{just status},
read-only in the same sandbox, reports the core release --- tag
\ReleaseTag, \TotalDecls\ declarations, \KernelTheorems\ proven,
\StatedCount\ stated, all gates \CertifiedStatus, tree clean, and a
93-artifact ledger --- as already certified from prior, checked-in
release state; nothing in this session rebuilt or re-derived that
standing, and nothing in this sandbox's missing toolchain retroactively
unproves a theorem the pinned toolchain already admitted at
certification time. STANDING.md is the governing record of that
certification, independent of any later environment's ability to replay
it locally. The limitation reported here is a property of this
particular writing environment's bootstrap state, worth recording
honestly precisely because this thesis's own discipline forbids asserting
replay success without having actually run it --- not a property, claimed
or implied, of the release the rest of this thesis reports on.
